\section{Limitations and Future Work}
\label{sec:limitations}

\subsection{Cold Start Problem}

The cold start problem is a fundamental limitation of collaborative filtering:

\textbf{New User Cold Start:}
A new user with no history cannot be modeled using collaborative filtering. The system falls back to 
the global mean rating, which is suboptimal. Real-world solutions include:

\begin{itemize}
  \item \textit{Content-based filtering}: Use item metadata (genre, cast) to make initial recommendations.
  
  \item \textit{Hybrid approaches}: Combine collaborative and content-based methods.
  
  \item \textit{Active learning}: Ask new users to rate items to quickly accumulate history.
  
  \item \textit{Knowledge graph integration}: Leverage external knowledge (social graphs, user profiles).
\end{itemize}

\textbf{New Item Cold Start:}
Similarly, newly released items have no rating history. Solutions include:

\begin{itemize}
  \item \textit{Release boost}: Temporarily promote new items to accumulate ratings faster.
  
  \item \textit{Content-based similarity}: Recommend new items similar to items the user enjoyed.
  
  \item \textit{Temporal modeling}: Track how ratings evolve over time and extrapolate for new items.
\end{itemize}

\subsection{Data Sparsity}

The rating matrix is 93.7\% sparse. This has implications:

\begin{enumerate}
  \item \textbf{Limited training signal}: Many users and items have few ratings, limiting factor estimation.
  
  \item \textbf{Genre and temporal bias}: Popular items and recent ratings dominate the training set, 
  potentially biasing predictions.
  
  \item \textbf{Extrapolation uncertainty}: Predictions for users/items with sparse data are less reliable.
\end{enumerate}

\textbf{Mitigation strategies:}
\begin{itemize}
  \item Data augmentation (synthetic ratings based on external signals).
  \item Ensemble methods combining multiple models.
  \item Bayesian approaches quantifying uncertainty.
\end{itemize}

\subsection{Scalability Constraints}

The current system has scalability limitations:

\begin{enumerate}
  \item \textbf{Model training}: SVD training is $O(k \times |\mathcal{R}| \times T)$, where $T$ is the number of epochs. 
  On large datasets (billions of ratings), this becomes expensive.
  
  \item \textbf{Inference latency}: Each prediction requires a dense matrix multiplication, which is fast but 
  becomes a bottleneck under extreme load (millions of predictions per second).
  
  \item \textbf{Model size}: Storing factors for all users and items requires memory proportional to 
  $(|\mathcal{U}| + |\mathcal{I}|) \times k$. For billion-scale datasets, this becomes prohibitive.
\end{enumerate}

\textbf{Solutions for large scale:}
\begin{itemize}
  \item Distributed training (PySpark, Ray, Horovod).
  \item Approximate inference (learned indexes, hashing).
  \item Model compression (quantization, pruning).
  \item Hierarchical approaches (cluster users/items first, then model within clusters).
\end{itemize}

\subsection{Model Stability and Concept Drift}

The current system assumes the model is trained once and deployed indefinitely. In practice:

\begin{enumerate}
  \item \textbf{Concept drift}: User preferences and item popularity change over time. A model trained 
  on historical data becomes stale.
  
  \item \textbf{New users/items}: Continuously arriving new users and items are not reflected in the 
  static model.
  
  \item \textbf{Feedback loops}: Recommendations influence user behavior, which in turn influences 
  future recommendations (potential bias amplification).
\end{enumerate}

\textbf{Production strategies:}
\begin{itemize}
  \item \textit{Periodic retraining}: Retrain the model weekly or daily on fresh data.
  
  \item \textit{Online learning}: Incrementally update the model with new ratings.
  
  \item \textit{A/B testing}: Test new models on subsets of users before full deployment.
  
  \item \textit{Monitoring}: Track prediction quality metrics and alert on degradation.
\end{itemize}

\subsection{Fairness and Bias}

Recommendation systems can perpetuate or amplify biases:

\begin{enumerate}
  \item \textbf{Popularity bias}: Popular items receive more ratings and are more visible in recommendations, 
  creating a rich-get-richer dynamic.
  
  \item \textbf{Filter bubbles}: Recommending items similar to a user's history can reinforce existing 
  preferences and limit diversity.
  
  \item \textbf{Demographic bias}: If the training data reflects demographic stereotypes, recommendations 
  will propagate these biases.
  
  \item \textbf{Cold-start amplification}: New items and new users may be underrepresented due to cold-start 
  effects, creating systemic disadvantage.
\end{enumerate}

\textbf{Mitigation approaches:}
\begin{itemize}
  \item \textit{Diversity objectives}: Optimize for diversity in addition to accuracy.
  
  \item \textit{Fair allocation}: Quota-based recommendations to ensure representation.
  
  \item \textit{Explainability}: Help users understand why items were recommended.
  
  \item \textit{User control}: Allow users to adjust recommendation preferences and opt-out.
  
  \item \textit{Audit and monitoring}: Regularly audit recommendations for bias.
\end{itemize}

This system does not explicitly address fairness; it is an important consideration for production systems.

\subsection{Interpretability and Explainability}

While SVD factors are more interpretable than neural network embeddings, they remain opaque:

\begin{enumerate}
  \item \textbf{Factor semantics}: The $k$ latent factors represent dimensions of preference space, but 
  their meaning is not labeled or easily understood.
  
  \item \textbf{Prediction explanation}: Why did the model predict 4.0 for ``user 196, movie 242''? 
  The answer involves complex factor combinations.
  
  \item \textbf{Error attribution}: When predictions are wrong, it's difficult to diagnose the cause 
  (sparse data? model overfitting? concept drift?).
\end{enumerate}

\textbf{Solutions:}
\begin{itemize}
  \item \textit{Post-hoc explanations}: Factor analysis, nearest neighbor look-ups, LIME/SHAP.
  
  \item \textit{Attention mechanisms}: Newer models provide interpretable attention weights.
  
  \item \textit{User interviews}: Qualitative feedback from users about recommendation quality.
\end{itemize}

\subsection{Privacy Considerations}

Recommendation systems may raise privacy concerns:

\begin{enumerate}
  \item \textbf{User tracking}: Collecting and storing user ratings reveals preferences.
  
  \item \textbf{Model inversion}: Sophisticated attacks may infer user data from models.
  
  \item \textbf{Data retention}: Long-term storage of user behavior creates security risks.
  
  \item \textbf{Regulatory compliance}: GDPR, CCPA, and other regulations impose requirements on 
  data handling and user rights.
\end{enumerate}

\textbf{Practices:}
\begin{itemize}
  \item \textit{Data minimization}: Collect only necessary data.
  
  \item \textit{Anonymization}: Hash user IDs and remove identifying information.
  
  \item \textit{Federated learning}: Train models on decentralized data without collecting it centrally.
  
  \item \textit{Right to explanation}: Provide users with explanations of recommendations.
  
  \item \textit{Right to deletion}: Allow users to request deletion of their data.
\end{itemize}

\subsection{Future Research Directions}

Promising directions for extending this system:

\begin{enumerate}
  \item \textbf{Deep learning}: Implement neural collaborative filtering (NCF, autoencoders) for improved 
  expressiveness.
  
  \item \textbf{Context-aware recommendations}: Incorporate context (time, location, device) into predictions.
  
  \item \textbf{Explainable AI}: Develop explainability methods tailored to recommendations.
  
  \item \textbf{Fairness-aware recommendations}: Optimize jointly for accuracy and fairness.
  
  \item \textbf{Real-time learning}: Implement online learning for immediate model updates.
  
  \item \textbf{Metamodeling}: Learn to select and combine multiple models dynamically.
  
  \item \textbf{Federated recommendations}: Enable privacy-preserving distributed training.
\end{enumerate}

Despite these limitations, the current system serves as a robust baseline and production-ready implementation 
for movie rating prediction, suitable for educational and experimental purposes.
